\documentclass[10pt,a4paper]{amsart}
\usepackage{amsmath,amssymb,amsthm,graphicx}
\usepackage{verbatim}
\newenvironment{reference}{\comment}{\endcomment}
%\newenvironment{reference}{}{}
\newenvironment{slogan}{\comment}{\endcomment}
\newenvironment{history}{\comment}{\endcomment}

% For commutative diagrams we use Xy-pic
\usepackage[all]{xy}

% We use 2cell for 2-commutative diagrams.
\xyoption{2cell}
% Theorem environments.
%
\theoremstyle{plain}
\newtheorem{theorem}[subsection]{Theorem}
\newtheorem{proposition}[subsection]{Proposition}
\newtheorem{lemma}[subsection]{Lemma}

\theoremstyle{definition}
\newtheorem{definition}[subsection]{Definition}
\newtheorem{example}[subsection]{Example}
\newtheorem{exercise}[subsection]{Exercise}
\newtheorem{situation}[subsection]{Situation}

\theoremstyle{remark}
\newtheorem{remark}[subsection]{Remark}
\newtheorem{remarks}[subsection]{Remarks}

\numberwithin{equation}{subsection}

% Macros
%
\def\lim{\mathop{\mathrm{lim}}\nolimits}
\def\colim{\mathop{\mathrm{colim}}\nolimits}
\def\Spec{\mathop{\mathrm{Spec}}}
\def\Hom{\mathop{\mathrm{Hom}}\nolimits}
\def\Ext{\mathop{\mathrm{Ext}}\nolimits}
\def\SheafHom{\mathop{\mathcal{H}\!\mathit{om}}\nolimits}
\def\SheafExt{\mathop{\mathcal{E}\!\mathit{xt}}\nolimits}
\def\Sch{\mathit{Sch}}
\def\Mor{\mathop{\mathrm{Mor}}\nolimits}
\def\Ob{\mathop{\mathrm{Ob}}\nolimits}
\def\Sh{\mathop{\mathit{Sh}}\nolimits}
\def\NL{\mathop{N\!L}\nolimits}
\def\CH{\mathop{\mathrm{CH}}\nolimits}
\def\proetale{{pro\text{-}\acute{e}tale}}
\def\etale{{\acute{e}tale}}
\def\QCoh{\mathit{QCoh}}
\def\Ker{\mathop{\mathrm{Ker}}}
\def\Im{\mathop{\mathrm{Im}}}
\def\Coker{\mathop{\mathrm{Coker}}}
\def\Coim{\mathop{\mathrm{Coim}}}

% Boxtimes
%
\DeclareMathSymbol{\boxtimes}{\mathbin}{AMSa}{"02}

%
% Macros for moduli stacks/spaces
%
\def\QCohstack{\mathcal{QC}\!\mathit{oh}}
\def\Cohstack{\mathcal{C}\!\mathit{oh}}
\def\Spacesstack{\mathcal{S}\!\mathit{paces}}
\def\Quotfunctor{\mathrm{Quot}}
\def\Hilbfunctor{\mathrm{Hilb}}
\def\Curvesstack{\mathcal{C}\!\mathit{urves}}
\def\Polarizedstack{\mathcal{P}\!\mathit{olarized}}
\def\Complexesstack{\mathcal{C}\!\mathit{omplexes}}
% \Pic is the operator that assigns to X its picard group, usage \Pic(X)
% \Picardstack_{X/B} denotes the Picard stack of X over B
% \Picardfunctor_{X/B} denotes the Picard functor of X over B
\def\Pic{\mathop{\mathrm{Pic}}\nolimits}
\def\Picardstack{\mathcal{P}\!\mathit{ic}}
\def\Picardfunctor{\mathrm{Pic}}
\def\Deformationcategory{\mathcal{D}\!\mathit{ef}}
\usepackage{cleveref}
\crefname{originalequation}{original equation}{original equations}
\Crefname{originalequation}{Original equation}{Original equations}
\makeatletter
\newlabel{equation-compare-with-pullback}{{0A9S}{source}{}{sourcepointer-1}{}}
\newlabel{equation-compare-with-pullback@cref}{{[originalequation][0][]0A9S}{[1][0][]source}{}{}{}}
\makeatother
\crefname{originallemma}{original lemma}{original lemmas}
\Crefname{originallemma}{Original lemma}{Original lemmas}
\makeatletter
\newlabel{lemma-proper-flat}{{0E4I}{source}{}{sourcepointer-2}{}}
\newlabel{lemma-proper-flat@cref}{{[originallemma][0][]0E4I}{[1][0][]source}{}{}{}}
\makeatother
\crefname{originaltheorem}{original theorem}{original theorems}
\Crefname{originaltheorem}{Original theorem}{Original theorems}
\makeatletter
\newlabel{perfect-theorem-DQCoh-is-Ddga}{{09M5}{source}{}{sourcepointer-3}{}}
\newlabel{perfect-theorem-DQCoh-is-Ddga@cref}{{[originaltheorem][0][]09M5}{[1][0][]source}{}{}{}}
\makeatother
\crefname{originaltheorem}{original theorem}{original theorems}
\Crefname{originaltheorem}{Original theorem}{Original theorems}
\makeatletter
\newlabel{perfect-theorem-bondal-van-den-Bergh}{{09IS}{source}{}{sourcepointer-4}{}}
\newlabel{perfect-theorem-bondal-van-den-Bergh@cref}{{[originaltheorem][0][]09IS}{[1][0][]source}{}{}{}}
\makeatother
\crefname{originalreference}{original reference}{original references}
\Crefname{originalreference}{Original reference}{Original references}
\makeatletter
\newlabel{more-algebra-remark-P-resolution}{{09PB}{source}{}{sourcepointer-5}{}}
\newlabel{more-algebra-remark-P-resolution@cref}{{[originalreference][0][]09PB}{[1][0][]source}{}{}{}}
\makeatother
\crefname{originallemma}{original lemma}{original lemmas}
\Crefname{originallemma}{Original lemma}{Original lemmas}
\makeatletter
\newlabel{lemma-compare-with-pullback-perfect}{{0A9T}{source}{}{sourcepointer-6}{}}
\newlabel{lemma-compare-with-pullback-perfect@cref}{{[originallemma][0][]0A9T}{[1][0][]source}{}{}{}}
\makeatother
\title{The comparison map between pullback and the right adjoint is an isomorphism for a proper flat finite presentation morphism}
\author{The Stacks Project Authors}
\date{stacks-project:duality.tex:local-snapshot-20261012; GNU FDL 1.2 or later}
\begin{document}
\maketitle
\noindent\textit{Editorial note.} All delivered mathematical text is the author's own, copied from the pinned Stacks Project chapter file duality.tex; no statement or proof step was rewritten, reordered or removed.\par\medskip
\noindent\textit{Editorial note.} Preamble provenance: the pinned chapter file opens with an input of the shared preamble and carries no preamble text of its own, so the project's single Stacks-wide preamble staged at sources/stacks-duality/preamble.tex is excerpted verbatim under the recipe fixed by LANE\_BRIEF 22.1, as three spans, lines 8-13 (verbatim package and the reference/slogan/history environments), lines 14-18 (Xy-pic setup) and lines 154-217 (the native theorem declarations and mathematical macros). Nothing was added, removed or redefined; the wrapper class and the amsmath/amssymb/amsthm/graphicx base packages are supplied by the assembler.\par\medskip
\noindent\textit{Editorial note.} The counted statement and proof contain no cross reference, no citation and no included file, so the delivered body is self contained.\par\medskip
\noindent\textit{Editorial note.} Canonical location of the counted proposition: Stacks tag 0E4K (sources/stacks-duality/pinned-original-source.tex line 2892).\par\medskip
\noindent\textit{Allowed standard prerequisites (editorial).}\begin{itemize}
\item The right adjoint commutes with direct sums for a proper flat morphism. Same chapter, Lemma labelled lemma-proper-flat, cited at proof lines 2919-2921. Makes the isomorphism locus closed under direct sums.
\item Generation of the quasi-coherent derived category by a single perfect object, in a module category. Derived Categories of Schemes, Theorems labelled perfect-theorem-DQCoh-is-Ddga and perfect-theorem-bondal-van-den-Bergh, cited at proof lines 2927-2933, and More on Algebra, Remark labelled more-algebra-remark-P-resolution, cited at proof lines 2931-2933. Reduces the claim to one perfect object.
\item The comparison map is an isomorphism for perfect objects. Same chapter, Lemma labelled lemma-compare-with-pullback-perfect, cited at proof lines 2933-2935. Settles the reduced case.
\end{itemize}\medskip
\ifdefined\hypertarget\hypertarget{sourcepointer-1}{}\fi
\noindent\textit{Source reference.} Original equation 0A9S: sources/stacks-duality/pinned-original-source.tex, source lines 2079--2079. Same-chapter pointer into the author environment equation at source line 2079; invoked at proof line 2903 of the counted proof. It is retained only as a navigation pointer; every load-bearing step is inside the delivered span.\par\medskip
\ifdefined\hypertarget\hypertarget{sourcepointer-2}{}\fi
\noindent\textit{Source reference.} Original lemma 0E4I: sources/stacks-duality/pinned-original-source.tex, source lines 2831--2831. Same-chapter pointer into the author environment lemma at source line 2831; invoked at proof line 2901 of the counted proof. It is retained only as a navigation pointer; every load-bearing step is inside the delivered span.\par\medskip
\ifdefined\hypertarget\hypertarget{sourcepointer-3}{}\fi
\noindent\textit{Source reference.} Original theorem 09M5: sources/stacks-chow/tags.txt, source lines 12168--12168. Cross-chapter standard prerequisite located in the project tag list staged in the local snapshot; invoked at proof line 2908 of the counted proof. It is retained only as a navigation pointer; every load-bearing step is inside the delivered span.\par\medskip
\ifdefined\hypertarget\hypertarget{sourcepointer-4}{}\fi
\noindent\textit{Source reference.} Original theorem 09IS: sources/stacks-chow/tags.txt, source lines 12050--12050. Cross-chapter standard prerequisite located in the project tag list staged in the local snapshot; invoked at proof line 2910 of the counted proof. It is retained only as a navigation pointer; every load-bearing step is inside the delivered span.\par\medskip
\ifdefined\hypertarget\hypertarget{sourcepointer-5}{}\fi
\noindent\textit{Source reference.} Original reference 09PB: sources/stacks-chow/tags.txt, source lines 12244--12244. Cross-chapter standard prerequisite located in the project tag list staged in the local snapshot; invoked at proof line 2912 of the counted proof. It is retained only as a navigation pointer; every load-bearing step is inside the delivered span.\par\medskip
\ifdefined\hypertarget\hypertarget{sourcepointer-6}{}\fi
\noindent\textit{Source reference.} Original lemma 0A9T: sources/stacks-duality/pinned-original-source.tex, source lines 2085--2085. Same-chapter pointer into the author environment lemma at source line 2085; invoked at proof line 2916 of the counted proof. It is retained only as a navigation pointer; every load-bearing step is inside the delivered span.\par\medskip
\begin{lemma}
\label{lemma-compare-with-pullback-flat-proper}
Let $Y$ be a quasi-compact and quasi-separated scheme.
Let $f : X \to Y$ be a morphism of schemes which is proper, flat, and
of finite presentation.
The map (\ref{equation-compare-with-pullback}) is an isomorphism
for every object $K$ of $D_\QCoh(\mathcal{O}_Y)$.
\end{lemma}

\begin{proof}
By Lemma \ref{lemma-proper-flat} we know that $a$ commutes
with direct sums. Hence the collection of objects of
$D_\QCoh(\mathcal{O}_Y)$ for which (\ref{equation-compare-with-pullback})
is an isomorphism is a strictly full, saturated, triangulated
subcategory of $D_\QCoh(\mathcal{O}_Y)$ which is moreover
preserved under taking direct sums. Since $D_\QCoh(\mathcal{O}_Y)$
is a module category (Derived Categories of Schemes, Theorem
\ref{perfect-theorem-DQCoh-is-Ddga}) generated by a single
perfect object (Derived Categories of Schemes, Theorem
\ref{perfect-theorem-bondal-van-den-Bergh})
we can argue as in
More on Algebra, Remark \ref{more-algebra-remark-P-resolution}
to see that it suffices to prove (\ref{equation-compare-with-pullback})
is an isomorphism for a single perfect object.
However, the result holds for perfect objects, see
Lemma \ref{lemma-compare-with-pullback-perfect}.
\end{proof}

\end{document}
